% GENERATED FILE. Edit canonical lessons, then run npm run generate:books.
% canonical-source-hash: fnv1a64:afad04aba7fef6d0
% canonical-lessons: RU-C424-parom, RU-C424-paluba, RU-C424-kayuta, RU-C424-prichal, RU-C424-pristan

\chapter{On the ferry}
\label{ch:ru-on-the-ferry}
\hlchaptermodality{424}

\begin{chapteropening}
\textbf{By the end of this chapter:} \emph{I can travel by ferry: the ferry, the deck, the cabin, the mooring and the landing stage.}

Complete the last lesson of chapter 424: I can travel by ferry: the ferry, the deck, the cabin, the mooring and the landing stage.
\end{chapteropening}

\section[\texorpdfstring{\ru{паром} (paróm)}{паром (paróm)}]{\ru{паром} (paróm) — a ferry}

\label{lesson:RU-C424-parom}



\begin{quote}
Before the new one: say the Russian for to cancel, then the Russian for compensation.
\end{quote}

\subsection*{You'll want to know: \ru{паром}}
\textbf{\ru{паром}} (\emph{paróm}) — \textquotedblleft{}a ferry\textquotedblright{}.

\begin{quote}
\ru{На} \ru{остров} \ru{ходит} \ru{паром}.
\end{quote}

\begin{quote}
\emph{Na óstrov khódit paróm.}
\end{quote}

\begin{quote}
A ferry goes to the island.
\end{quote}

Masculine. \emph{\ru{На} \ru{пароме}} — by ferry.

\begin{cousinweb}
\emph{\ru{паром}} is inherited from Proto-Slavic \emph{pormŭ} 'a ferry' (EtymDB).
\end{cousinweb}

\begin{grammarlens}[title={What you've built}]
One more word for telling, explaining and describing at B1.
\end{grammarlens}

\subsection*{Your turn}
\begin{itemize}
\raggedright
  \item \emph{Say it:} \emph{paróm}
  \item \emph{Say it:} \emph{paróm}, once more
  \item \emph{Say it:} a sentence of your own with \emph{paróm}
  \item \emph{Recall:} say the Russian for to cancel, then the Russian for compensation, then say \emph{paróm} again
\end{itemize}

\begin{tcolorbox}[breakable,colback=teal!4,colframe=teal!35!black,title={Before you move on}]
What does \ru{паром} mean? (\textquotedblleft{}A ferry\textquotedblright{}.) Say it once more.
\end{tcolorbox}
\section[\texorpdfstring{\ru{палуба} (páluba)}{палуба (páluba)}]{\ru{палуба} (páluba) — a deck (of a ship)}

\label{lesson:RU-C424-paluba}



\begin{quote}
Before the new one: say the Russian for compensation, then the Russian for a ferry.
\end{quote}

\subsection*{You'll want to know: \ru{палуба}}
\textbf{\ru{палуба}} (\emph{páluba}) — \textquotedblleft{}a deck (of a ship)\textquotedblright{}.

\begin{quote}
\ru{Вечером} \ru{все} \ru{пассажиры} \ru{сидели} \ru{на} \ru{палубе}.
\end{quote}

\begin{quote}
\emph{Vécherom vse passazhíry sidéli na pálube.}
\end{quote}

\begin{quote}
In the evening all the passengers sat on deck.
\end{quote}

Feminine; the stress is on the first syllable.

\begin{cousinweb}
The origin of \emph{\ru{палуба}} is not settled; EtymDB lists the word without a source.
\end{cousinweb}

\begin{grammarlens}[title={What you've built}]
One more word for telling, explaining and describing at B1.
\end{grammarlens}

\subsection*{Your turn}
\begin{itemize}
\raggedright
  \item \emph{Say it:} \emph{páluba}
  \item \emph{Say it:} \emph{páluba}, once more
  \item \emph{Say it:} a sentence of your own with \emph{páluba}
  \item \emph{Recall:} say the Russian for compensation, then the Russian for a ferry, then say \emph{páluba} again
\end{itemize}

\begin{tcolorbox}[breakable,colback=teal!4,colframe=teal!35!black,title={Before you move on}]
What does \ru{палуба} mean? (\textquotedblleft{}A deck (of a ship)\textquotedblright{}.) Say it once more.
\end{tcolorbox}
\section[\texorpdfstring{\ru{каюта} (kayúta)}{каюта (kayúta)}]{\ru{каюта} (kayúta) — a cabin (on a ship)}

\label{lesson:RU-C424-kayuta}



\begin{quote}
Before the new one: say the Russian for a ferry, then the Russian for a deck.
\end{quote}

\subsection*{You'll want to know: \ru{каюта}}
\textbf{\ru{каюта}} (\emph{kayúta}) — \textquotedblleft{}a cabin (on a ship)\textquotedblright{}.

\begin{quote}
\ru{У} \ru{нас} \ru{была} \ru{маленькая} \ru{каюта} \ru{без} \ru{окна}.
\end{quote}

\begin{quote}
\emph{U nas bylá málenkaya kayúta bez okná.}
\end{quote}

\begin{quote}
We had a small cabin with no window.
\end{quote}

Feminine.

\begin{cousinweb}
\emph{\ru{каюта}} was borrowed from German \emph{Kajüte}, which EtymDB traces to Middle Low German \emph{kajüte}.
\end{cousinweb}

\begin{grammarlens}[title={What you've built}]
One more word for telling, explaining and describing at B1.
\end{grammarlens}

\subsection*{Your turn}
\begin{itemize}
\raggedright
  \item \emph{Say it:} \emph{kayúta}
  \item \emph{Say it:} \emph{kayúta}, once more
  \item \emph{Say it:} a sentence of your own with \emph{kayúta}
  \item \emph{Recall:} say the Russian for a ferry, then the Russian for a deck, then say \emph{kayúta} again
\end{itemize}

\begin{tcolorbox}[breakable,colback=teal!4,colframe=teal!35!black,title={Before you move on}]
What does \ru{каюта} mean? (\textquotedblleft{}A cabin (on a ship)\textquotedblright{}.) Say it once more.
\end{tcolorbox}
\section[\texorpdfstring{\ru{причал} (prichál)}{причал (prichál)}]{\ru{причал} (prichál) — a mooring, a berth}

\label{lesson:RU-C424-prichal}



\begin{quote}
Before the new one: say the Russian for a deck, then the Russian for a cabin.
\end{quote}

\subsection*{You'll want to know: \ru{причал}}
\textbf{\ru{причал}} (\emph{prichál}) — \textquotedblleft{}a mooring, a berth\textquotedblright{}.

\begin{quote}
\ru{Лодки} \ru{стоят} \ru{у} \ru{причала}.
\end{quote}

\begin{quote}
\emph{Lódki stoyát u prichála.}
\end{quote}

\begin{quote}
The boats are at the mooring.
\end{quote}

Masculine. Built on \emph{\ru{причалить}} — to moor.

\begin{cousinweb}
\emph{\ru{причал}} is built on \emph{\ru{причалить}} 'to moor', from \emph{\ru{при}-} and \emph{\ru{чалить}} 'to tie up (a boat)'. EtymDB lists \emph{\ru{чалить}} without a source, so its origin is left open.
\end{cousinweb}

\begin{grammarlens}[title={What you've built}]
One more word for telling, explaining and describing at B1.
\end{grammarlens}

\subsection*{Your turn}
\begin{itemize}
\raggedright
  \item \emph{Say it:} \emph{prichál}
  \item \emph{Say it:} \emph{prichál}, once more
  \item \emph{Say it:} a sentence of your own with \emph{prichál}
  \item \emph{Recall:} say the Russian for a deck, then the Russian for a cabin, then say \emph{prichál} again
\end{itemize}

\begin{tcolorbox}[breakable,colback=teal!4,colframe=teal!35!black,title={Before you move on}]
What does \ru{причал} mean? (\textquotedblleft{}A mooring, a berth\textquotedblright{}.) Say it once more.
\end{tcolorbox}
\section[\texorpdfstring{\ru{пристань} (prístan)}{пристань (prístan)}]{\ru{пристань} (prístan) — a landing stage, a jetty}

\label{lesson:RU-C424-pristan}



\begin{quote}
Before the new one: say the Russian for a cabin, then the Russian for a mooring.
\end{quote}

\subsection*{You'll want to know: \ru{пристань}}
\textbf{\ru{пристань}} (\emph{prístan}) — \textquotedblleft{}a landing stage, a jetty\textquotedblright{}.

\begin{quote}
\ru{Мы} \ru{ждали} \ru{паром} \ru{на} \ru{пристани}.
\end{quote}

\begin{quote}
\emph{My zhdáli paróm na prístani.}
\end{quote}

\begin{quote}
We waited for the ferry at the landing stage.
\end{quote}

Feminine, though it ends in -\ru{ь}: \emph{\ru{на} \ru{пристани}}.

\begin{cousinweb}
\emph{\ru{пристань}} is built on \emph{\ru{пристать}} 'to put in to shore', from \emph{\ru{при}-} and \emph{\ru{стать}} 'to stand, to come to a stop'. EtymDB lists \emph{\ru{пристать}} without a deeper source.
\end{cousinweb}

\begin{grammarlens}[title={What you've built}]
One more word for telling, explaining and describing at B1.
\end{grammarlens}

\subsection*{Your turn}
\begin{itemize}
\raggedright
  \item \emph{Say it:} \emph{prístan}
  \item \emph{Say it:} \emph{prístan}, once more
  \item \emph{Say it:} a sentence of your own with \emph{prístan}
  \item \emph{Recall:} say the Russian for a cabin, then the Russian for a mooring, then say \emph{prístan} again
\end{itemize}

\begin{tcolorbox}[breakable,colback=teal!4,colframe=teal!35!black,title={Before you move on}]
What does \ru{пристань} mean? (\textquotedblleft{}A landing stage, a jetty\textquotedblright{}.) Say it once more.
\end{tcolorbox}
